% !TEX root = ../main.tex

\chapter{Stacks and local representability}
\label[chapter]{chap:Coherence_Descent_Local_Representability}

A moduli problem assigns families to parameter objects. For example, a
family of zeros of $f\colon U\to\R^r$ parametrized by $T$ is a map
$T\to U$ satisfying the equation. In derived geometry the equation
includes a specified homotopy to zero. Pulling the family back along
$T'\to T$ gives a family over $T'$.

This viewpoint defines a moduli problem before a geometric representing
object has been constructed. Descent expresses that families and their
isomorphisms can be glued. The main result of this chapter gives a way
to recognize when the resulting stack is a derived scheme: construct
an open cover by representable stacks.

This returns to the coherence problem posed in
\Cref{sec:Local_Zero_Loci_Not_Enough}. Starting with a globally defined
stack supplies the overlaps and their higher compatibilities; the work
is to construct open charts of that functor. We develop descent and
atlases with the vector bundle zero locus as a running example, then
formulate representability relative to a parameter stack. The distinction
between gluing families and gluing charts will be essential throughout.

\section{The open topology and descent}

Descent concerns a cover of a parameter object: it reconstructs an
anima of global families from coherent local families. We specify the
test objects and covers first, then express this reconstruction as a limit.

We use the category $\mathrm{DMfd}$ of derived manifolds as test objects.
By the affine comparison, these are spectra of finitely presented animated
smooth rings. An open embedding is an inclusion of an open subobject
with its restricted structure sheaf. An \emph{open covering family}
$\{T_i\to T\}$ is a family of open embeddings whose underlying images
cover $|T|$. Smooth localization describes these embeddings algebraically.
They are preserved by pullback.

\begin{definition}[Prestack]
\label[definition]{def:Prestack_Derived_Cinfty}
A derived smooth prestack is a functor
$\mathcal F\colon\mathrm{DMfd}^{\mathrm{op}}\to\An$.
\end{definition}

The anima $\mathcal F(T)$ contains the $T$-families, their isomorphisms,
and higher identifications between those isomorphisms.
Replacing it by its set of components loses automorphisms as well as
higher coherence. Pullback along a map of parameter objects is part
of the functor.

For an open cover, write $T_{i_0\cdots i_n}$ for the iterated intersection.
Restriction gives a cosimplicial diagram with degree-$n$ term
\[
 C^n=\prod_{i_0,\ldots,i_n}\mathcal F(T_{i_0\cdots i_n}).
\]
Its face maps restrict to smaller intersections, and its degeneracy maps
repeat indices. The following formula for its limit avoids any
need for an infinite disjoint union to be a test object.

\begin{definition}[Stack]
\label[definition]{def:Stack_Derived_Cinfty}
A prestack is a stack if, for every open covering family, the natural map
\begin{equation}\label{eq:Stack_Descent_Parameter_Cech}
 \mathcal F(T)\longrightarrow
 \lim_{[n]\in\Delta}
 \prod_{i_0,\ldots,i_n}\mathcal F(T_{i_0\cdots i_n})
\end{equation}
is an isomorphism of animae.
\end{definition}

For a set-valued functor this reduces to the usual sheaf condition:
local elements must agree on double overlaps.
For groupoid-valued functors one instead chooses isomorphisms on double
overlaps and imposes the cocycle condition on triple overlaps.
For general animae the cocycle condition itself is supplied by a homotopy,
followed by higher compatibility data. The limit in
\Cref{eq:Stack_Descent_Parameter_Cech} packages all of it.

Descent applies to maps between families too. It says that the anima of
global families is the anima of coherent local data, rather than merely
that some global family exists.

The resulting $\infty$-category of stacks is denoted
$\mathrm{dSt}_{\Cinfty}^{\mathrm{fp}}$ here.
The superscript records our site of finitely presented test objects.
Steffens embeds this category fully faithfully into his category using
all affine derived smooth schemes, and calls its objects stacks locally
of finite presentation
\cite[Proposition~2.1.40 and Definition~2.1.41]
{Steffens_Representability_Elliptic_Moduli}.
Our constructions take place in this subcategory.

\section{Represented stacks and the zero-locus example}

Representability asks whether all families arise as maps to a geometric
object. The universal property of a derived pullback answers this
question for the zero-locus functor and gives a model for the later
elliptic moduli problem.

For a derived scheme $X$ locally of finite presentation, its functor
of points is
\[
 h_X(T)=\Hom(T,X).
\]
It satisfies descent because morphisms of locally ringed objects can
be glued over an open cover of their source.
This is the subcanonical property of the open topology.
Moreover, the functor $X\mapsto h_X$ is fully faithful.

\begin{definition}[Representability]
\label[definition]{def:Represented_Stack}
A stack $\mathcal F$ is represented by a derived $\Cinfty$-scheme $X$
if there is a natural isomorphism
$\mathcal F(T)\cong\Hom(T,X)$ for every test object $T$.
\end{definition}

The representing object is unique. If $X$ is not affine, the notation
$\mathcal F(X)$ means the value obtained by descent from an affine
open cover. Under representability, the identity of $X$ corresponds
to a universal family in $\mathcal F(X)$. Every other family is its
pullback, by the Yoneda lemma.

Let $f\colon U\to\R^r$ be smooth. Define
\begin{equation}\label{eq:Derived_Zero_Locus_Prestack}
 \mathcal Z_f(T)=\Hom(T,U)\times_{\Hom(T,\R^r)}\Hom(T,*),
\end{equation}
with pullback in animae.
A point is a map $u\colon T\to U$ together with a homotopy from $f\circ u$
to zero. Identifications in this pullback give the isomorphisms between
such data and their higher compatibilities. Limits of stacks are computed on each test object, so
$\mathcal Z_f$ is a stack.

By the universal property of a pullback,
\[
 \mathcal Z_f(T)\cong\Hom(T,U\times_{\R^r}^{\mathrm{der}}*).
\]
Thus it is represented by the derived zero locus computed in \Cref{chap:Derived_Zero_Loci_And_Tangent_Complexes}.
This is a simple instance of the representability question that will
later arise for differential equations.

For a section $s\colon U\to E$ of a vector bundle, the analogous formula
is most conveniently written in the category over $U$:
\[
 \mathcal Z_s=U\times_E^{\mathrm{der}}U,
\]
using $s$ and the zero section. Restricting to a trivializing open
$U_i$ gives the Koszul presentation for its list of local components.
These local representatives will form an open atlas of $\mathcal Z_s$.

\section{Smooth stacks as derived stacks}

A smooth stack is a sheaf of animae on ordinary manifolds with their
open-cover topology. The inclusion of manifolds into derived manifolds
induces a fully faithful functor
\[
 j_!\colon\mathrm{Shv}(\mathrm{Mfd})
 \longrightarrow\mathrm{dSt}_{\Cinfty}^{\mathrm{fp}}.
\]
It is obtained by left Kan extension and sheafification, and sends
an ordinary manifold to its derived functor of points.
We can therefore regard smooth stacks as particular derived stacks.

The derived values of $j_!\mathcal F$ are not obtained just by
evaluating $\mathcal F$ on an underlying point set.
They are determined by this extension construction.
Conversely, restriction to ordinary test manifolds does not recover
every derived stack. For example, a multiple point and an ordinary
point have the same families over ordinary manifolds, while maps
from infinitesimal derived tests distinguish them.

There is also an essential distinction between the two categories of
pullbacks. In smooth stacks,
\[
 *\times_\R *=*
\]
for the two origins. After passing to derived stacks, the pullback
is the derived self-intersection, with algebra $\R[\varepsilon]$,
where $|\varepsilon|=1$ and $d\varepsilon=0$.
Thus $j_!$ does not preserve arbitrary pullbacks.
The pullbacks needed in the analytic argument will be compared using
submersiveness in \Cref{chap:Elliptic_Representability_Smooth_Derived_Bridge}.

\section{Open atlases and local representability}

We now turn from gluing families over a parameter cover to presenting
the moduli stack itself by charts. Here a \v{C}ech diagram reconstructs
the stack as a colimit. Effectiveness expresses that the charts cover
the functor; openness makes it possible to glue their representing
objects into a derived scheme.

Let $p\colon\mathcal U\to\mathcal F$ be a morphism of stacks.
Its \v{C}ech nerve is the simplicial stack
\[
 \mathcal U_n=
 \underbrace{\mathcal U\times_{\mathcal F}\cdots
 \times_{\mathcal F}\mathcal U}_{n+1\text{ factors}}.
\]
The augmentation maps its geometric realization to $\mathcal F$.
Here geometric realization means a colimit in stacks, computed by
taking the presheaf colimit and then sheafifying.

\begin{definition}[Effective epimorphism]
\label[definition]{def:Effective_Epimorphism}
The map $p$ is an effective epimorphism if
$|\mathcal U_\bullet|\to\mathcal F$ is an isomorphism.
\end{definition}

On the open site there is a useful test: every family in $\mathcal F(T)$
must, after an open cover of $T$, lift to $\mathcal U$, with an
identification of its image with the given family.
Equivalently, the induced map of sheaves of components is surjective.
This characterizes effective epimorphisms in the $\infty$-topos of stacks.

Starting with a globally defined stack supplies the overlap data.
If $U_i\to\mathcal F$ and $U_j\to\mathcal F$ are two charts, their
overlap is $U_i\times_{\mathcal F}U_j$. Its associativity and higher
compatibilities come from pullbacks in the ambient $\infty$-category.
The following condition makes these overlaps geometric open subsets.

\begin{definition}[Open atlas]
\label[definition]{def:Open_Atlas}
A morphism $\mathcal V\to\mathcal F$ is an open inclusion if every
pullback along $T\to\mathcal F$ is represented by an open subobject of $T$.
An open atlas is a family of open inclusions $h_{U_i}\to\mathcal F$
from derived schemes such that
$\coprod_i h_{U_i}\to\mathcal F$ is an effective epimorphism.
\end{definition}

The coproduct here is in stacks. Each map is a condition on families:
after choosing a family over $T$, belonging to the indicated chart
cuts out an open subset of $T$.

The word \emph{open} is crucial for representing a stack by a scheme.
For example, a nontrivial finite group $G$ has a classifying stack $BG$.
The map $*\to BG$ is an effective epimorphism, expressing local
triviality of principal $G$-bundles. It is not an open inclusion:
its pullback along itself is $G\to*$.
The automorphisms of the trivial bundle also show that $BG$ cannot be
represented by a derived scheme, whose real points form a set.
More general notions of atlas lead to orbifolds and other geometric
stacks, but an arbitrary effective epimorphism is insufficient for
scheme representability.

An open atlas contains exactly the data needed to reconstruct a
derived scheme from its local presentations.

\begin{theorem}[Open local representability]
\label[theorem]{thm:Open_Local_Representability}
A derived smooth stack is represented by a derived $\Cinfty$-scheme
locally of finite presentation if and only if it has an open atlas
by affine derived schemes of finite presentation.
One may equivalently allow derived schemes locally of finite presentation
as the atlas objects.
\end{theorem}

This is \citet[Proposition~2.1.48]{Steffens_Representability_Elliptic_Moduli}.

\begin{proof}[Gluing argument]
A represented stack has such an atlas by the definition of a derived
scheme. Conversely, set $U_{ij}=U_i\times_{\mathcal F}U_j$.
Since each chart map is open, $U_{ij}$ is represented by an open
subscheme of both $U_i$ and $U_j$.
The two maps identify these opens. On triple and higher overlaps
the identifications have the compatibility supplied by the \v{C}ech nerve.

Derived schemes glue along this open descent datum, producing a scheme
$X$. On the underlying topological level this is the usual gluing of
open subsets; the animated structure sheaves glue by descent.
The functor of points of $X$ realizes the same \v{C}ech diagram.
Since the atlas is an effective epimorphism,
$h_X\cong|\mathcal U_\bullet|\cong\mathcal F$.
Allowing non-affine charts gives the same criterion by refining each
one with an affine open cover.
\end{proof}

In the bundle-section example, the charts are the zero loci over a
trivializing cover of $U$. Their overlaps are the zero loci over the
corresponding intersections. The transition maps of the bundle identify
the Koszul models, and the theorem recovers the global derived zero locus.

The same argument preserves a local property such as quasi-smoothness.
If the charts are quasi-smooth, the cotangent amplitude condition
holds on the resulting scheme.

\section{Relative representability and base change}

Moduli problems often have a parameter stack $S$.
For example, the metric, the domain manifold, or a coefficient in an
equation may vary. The correct representability statement is then
relative to $S$.

\begin{definition}
\label[definition]{def:Relative_Representability}
A map $\mathcal F\to S$ is representable by derived $\Cinfty$-schemes
locally of finite presentation if, for each affine derived test
$T\to S$, the pullback $\mathcal F\times_S T$ is represented by such
a scheme. It is relatively quasi-smooth if these representing
morphisms to $T$ are quasi-smooth.
\end{definition}

This asserts representability of every pulled-back family.
It does not assert that the total stack $\mathcal F$ is itself a scheme
when $S$ has automorphisms.

\begin{corollary}[Locality on the base]
\label[corollary]{cor:Local_Representability_Relative}
Representability by derived schemes locally of finite presentation
can be checked after an effective epimorphic covering family of the base.
The same holds with the additional requirement of relative
quasi-smoothness.
\end{corollary}

\begin{proof}[Argument]
Pull the covering family back along a test object $T\to S$.
The local lifting property of an effective epimorphism gives an open
cover $\{T_j\to T\}$ on which these maps lift to members of the cover.
By assumption, the restrictions of $\mathcal F\times_S T$ to the
$T_j$ are represented. They form an open atlas of
$\mathcal F\times_S T$, so the theorem glues them.
Relative cotangent complexes commute with base change, and their
perfection and amplitude are local, giving the last assertion.
\end{proof}

The source formulation is
\cite[Definition~2.1.42 and Corollary~2.1.49]
{Steffens_Representability_Elliptic_Moduli}.
This permits a proof over a smooth stack to be carried out on manifold
charts of the base, provided the moduli construction is compatible
with pullback.

\section{Further reading}

The local representability criterion does not make the construction of
a chart automatic.
One must identify the entire derived functor on an open neighbourhood,
including its values on derived tests. An identification of real
solution sets, or even of their tangent complexes, proves less.
The later Kuranishi argument is designed to establish precisely
this functorial identification.

\begin{exercise}
Let $f\colon U\to\R^r$ and let $\{U_i\}$ cover $U$.
Use the functor-of-points formula to verify that the corresponding
zero loci form an effective open atlas of $\mathcal Z_f$.
\end{exercise}

\begin{exercise}
For a nontrivial finite group $G$, calculate $*\times_{BG}*$ and
explain both the effectiveness of $*\to BG$ and its failure to be
an open inclusion.
\end{exercise}

All the descent and representability results used in this chapter are
contained in \cite[Section~2.1]{Steffens_Representability_Elliptic_Moduli}.
